\documentclass[11pt,a4paper]{article}
\usepackage[margin=3cm]{geometry}
\usepackage{amsmath,amssymb,amsthm,mathrsfs}
\usepackage[T1]{fontenc}
\usepackage{microtype}
\usepackage[colorlinks=true,linkcolor=blue,citecolor=blue]{hyperref}

\theoremstyle{plain}
\newtheorem{theorem}{Theorem}[section]
\newtheorem{proposition}[theorem]{Proposition}
\newtheorem{lemma}[theorem]{Lemma}
\newtheorem{corollary}[theorem]{Corollary}
\theoremstyle{definition}
\newtheorem{remark}[theorem]{Remark}
\newtheorem{hypothesis}[theorem]{Hypothesis}

\newcommand{\A}{\mathfrak{A}}
\newcommand{\G}{\mathcal{G}}
\newcommand{\R}{\mathbb{R}}
\newcommand{\Z}{\mathbb{Z}}
\newcommand{\Om}{\Omega}
\newcommand{\ip}[2]{\langle #1,#2\rangle}
\DeclareMathOperator{\dist}{dist}
\DeclareMathOperator{\supp}{supp}
\DeclareMathOperator{\spec}{spec}
\DeclareMathOperator{\gapop}{gap}

\title{M3: clustering, the symmetry-breaking dichotomy,\\ and the reduction of uniqueness}
\author{}
\date{}

\begin{document}
\maketitle

\begin{abstract}
M3 asks for uniqueness of the ground state of $(\A,\alpha)$. This note delivers three things and
is explicit about what it does not deliver.

\smallskip
\noindent\textbf{(1) A correction.} Section~4.2 of the strategy note asserted that a uniform gap
gives exponential clustering and that ``uniqueness of a translation-invariant, exponentially
clustering ground state follows from a standard ergodicity argument''. The second half is
\emph{false}: clustering gives \emph{extremality}, not uniqueness --- two distinct pure phases both
cluster. The claim is withdrawn (\S\ref{sec:notuniq}).

\smallskip
\noindent\textbf{(2) A theorem, proved in full.} A uniform spectral gap plus the Glimm--Jaffe
\emph{exact} light cone gives exponential clustering with an absolute constant and the optimal
rate:
\[
  \bigl|\omega_g(AB)-\omega_g(A)\omega_g(B)\bigr|\ \le\ 2\|A\|\,\|B\|\,e^{-\gamma R/2},
  \qquad R=\dist(O_1,O_2),\ \ \gamma R\ge1 .
\]
The proof (\S\ref{sec:clust}) is self-contained: an exact spectral-support identity, one explicit
test function, and one optimisation. No Lieb--Robinson bound is needed, and no tail is paid,
because the light cone is sharp.

\smallskip
\noindent\textbf{(3) A dichotomy.} Every limit point of $(\omega_{\tilde g_n})$ is
$\Z_2$-invariant \cite[Lem.~3.3]{M0}. Combined with (2), this forces: \emph{if the $\Z_2$
symmetry is broken, then $\inf_n\gapop H(\tilde g_n)=0$}. So the uniform-gap hypothesis (H3) is not
a technical convenience --- it is equivalent to absence of symmetry breaking along the sequence
(\S\ref{sec:dich}).

\smallskip
\noindent\textbf{What is not delivered.} Uniqueness itself. \S\ref{sec:red} reduces it to four
precisely stated hypotheses (R1)--(R4) and proves the implications between them; (R3), a reverse
Osterwalder--Schrader reconstruction, is the one with no proof in the literature.
\textbf{M3 is therefore half-closed.}
\end{abstract}

\section{Setting}

Notation as in \cite{M0,M1}. Recall from \cite[Fact~1.3]{M0} the \emph{exact} light cone
$\alpha_t(\A(O))\subseteq\A(O_{|t|})$, and that $\alpha_t(A)=e^{itH(g)}Ae^{-itH(g)}$ for
$A\in\A(O)$ whenever $g\equiv1$ on $O_{|t|}$. Locality of the net gives $[\A(O),\A(O')]=0$ for
disjoint $O,O'$.

\begin{hypothesis}[uniform gap]\label{hyp:gap}
$\gamma:=\inf_n\gapop H(\tilde g_n)>0$, where $\gapop H(g)$ denotes the distance from $0$ to the
rest of $\spec H(g)$, positive for each $g$ by \cite[\S2.2--2.3]{GJII}.
\end{hypothesis}

\section{Exponential clustering from a gap and a sharp light cone}\label{sec:clust}

Throughout this section $g\in\G$ is fixed, $\gamma:=\gapop H(g)>0$,
$P_0:=|\Om_g\rangle\langle\Om_g|$, and $A\in\A(O_1)$, $B\in\A(O_2)$ with
$R:=\dist(O_1,O_2)>0$. We assume $d(g,O_2)>R$, so that the light cone representation of
$\alpha_t(B)$ is available for $|t|<R$.

\begin{lemma}[spectral supports and the locality identity]\label{lem:FG}
Define, for $t\in\R$,
\[
  F(t):=\ip{A^*\Om_g}{(e^{itH(g)}-P_0)B\Om_g},\qquad
  G(t):=\ip{B^*\Om_g}{(e^{-itH(g)}-P_0)A\Om_g}.
\]
Then, with $M_0:=\|A\|\|B\|$:
\begin{enumerate}
\item[\rm(i)] $F(0)=\omega_g(AB)-\omega_g(A)\omega_g(B)$;
\item[\rm(ii)] $F(t)=\int_{[\gamma,\infty)}e^{itE}d\rho(E)$ and
      $G(t)=\int_{(-\infty,-\gamma]}e^{itE}d\tilde\sigma(E)$ for complex measures with
      $|\rho|,|\tilde\sigma|\le M_0$;
\item[\rm(iii)] $F(t)=G(t)$ for all $|t|<R$.
\end{enumerate}
\end{lemma}

\begin{proof}
(i) is immediate from $P_0$ and $\|\Om_g\|=1$.

(ii) Put $\rho(\cdot):=\ip{A^*\Om_g}{P_{H(g)}(\cdot)(1-P_0)B\Om_g}$. Since
$\spec H(g)\subseteq\{0\}\cup[\gamma,\infty)$ and $(1-P_0)$ removes the eigenvalue $0$,
$\supp\rho\subseteq[\gamma,\infty)$, and
$|\rho|(\R)\le\|A^*\Om_g\|\,\|(1-P_0)B\Om_g\|\le M_0$. The statement for $G$ follows the same way
after reflecting $E\mapsto-E$.

(iii) Using $e^{-itH(g)}\Om_g=\Om_g$,
\[
  \ip{A^*\Om_g}{e^{itH(g)}B\Om_g}=\ip{\Om_g}{Ae^{itH(g)}Be^{-itH(g)}\Om_g}=\omega_g\bigl(A\alpha_t(B)\bigr),
\]
\[
  \ip{B^*\Om_g}{e^{-itH(g)}A\Om_g}=\ip{\Om_g}{e^{itH(g)}Be^{-itH(g)}A\Om_g}=\omega_g\bigl(\alpha_t(B)A\bigr),
\]
while both $P_0$ terms equal $\omega_g(A)\omega_g(B)$. Hence
$F(t)-G(t)=\omega_g([A,\alpha_t(B)])$. For $|t|<R$ the light cone gives
$\alpha_t(B)\in\A(O_{2,|t|})$ with $O_{2,|t|}\cap O_1=\emptyset$, so the commutator vanishes.
\end{proof}

\begin{lemma}[the test function]\label{lem:test}
For $\alpha>0$ and $\varepsilon>0$ set
$f_{\alpha,\varepsilon}(t):=\dfrac{1}{2\pi i}\,\dfrac{e^{-t^2/2\alpha}}{t-i\varepsilon}\in L^1(\R)$.
Then
\[
  \check f_{\alpha,\varepsilon}(E):=\int_\R f_{\alpha,\varepsilon}(t)e^{itE}\,dt
  =\int_0^\infty\sqrt{\tfrac{\alpha}{2\pi}}\,e^{-\alpha(E-E')^2/2}e^{-\varepsilon E'}\,dE',
\]
so that $0\le\check f_{\alpha,\varepsilon}\le1$ and, as $\varepsilon\downarrow0$,
$\check f_{\alpha,\varepsilon}(E)\to\Phi(\sqrt\alpha E)$, $\Phi$ the standard normal distribution
function. Moreover $\int_{|t|\ge R}|f_{\alpha,\varepsilon}|\,dt\le\frac{\alpha}{\pi R^2}e^{-R^2/2\alpha}$
uniformly in $\varepsilon$.
\end{lemma}

\begin{proof}
$\int e^{itE}\bigl(2\pi i(t-i\varepsilon)\bigr)^{-1}dt=e^{-\varepsilon E}\mathbf 1_{E>0}$ by
residues, and multiplication by $e^{-t^2/2\alpha}$ becomes convolution with
$\sqrt{\alpha/2\pi}\,e^{-\alpha E^2/2}$; this gives the displayed formula, whence
$0\le\check f_{\alpha,\varepsilon}\le1$ and the stated limit by dominated convergence. For the tail,
$|t-i\varepsilon|\ge|t|$ gives
$\int_{|t|\ge R}|f_{\alpha,\varepsilon}|\le\frac1{2\pi}\int_{|t|\ge R}\frac{e^{-t^2/2\alpha}}{|t|}dt
\le\frac1{\pi R}\int_R^\infty e^{-t^2/2\alpha}dt\le\frac{\alpha}{\pi R^2}e^{-R^2/2\alpha}$,
using $\int_R^\infty e^{-t^2/2\alpha}dt\le\frac{\alpha}{R}e^{-R^2/2\alpha}$.
\end{proof}

\begin{theorem}[exponential clustering]\label{thm:clust}
With the hypotheses of this section, if $\gamma R\ge1$ then
\[
  \bigl|\omega_g(AB)-\omega_g(A)\omega_g(B)\bigr|\ \le\ 2\,\|A\|\,\|B\|\;e^{-\gamma R/2}.
\]
\end{theorem}

\begin{proof}
Fix $\alpha,\varepsilon>0$. Since $f_{\alpha,\varepsilon}\in L^1$ and $|\rho|,|\tilde\sigma|<\infty$,
Fubini and Lemma~\ref{lem:FG}(ii) give
\[
  \int f_{\alpha,\varepsilon}F=\int_{[\gamma,\infty)}\check f_{\alpha,\varepsilon}\,d\rho,
  \qquad
  \int f_{\alpha,\varepsilon}G=\int_{(-\infty,-\gamma]}\check f_{\alpha,\varepsilon}\,d\tilde\sigma .
\]
By Lemma~\ref{lem:FG}(iii), $F-G$ vanishes on $(-R,R)$, and $|F|,|G|\le M_0$, so by
Lemma~\ref{lem:test}
\[
  \Bigl|\int f_{\alpha,\varepsilon}(F-G)\Bigr|
  =\Bigl|\int_{|t|\ge R} f_{\alpha,\varepsilon}(F-G)\Bigr|
  \le 2M_0\,\frac{\alpha}{\pi R^2}e^{-R^2/2\alpha}.
\]
Letting $\varepsilon\downarrow0$ and using dominated convergence
($0\le\check f_{\alpha,\varepsilon}\le1$),
\begin{equation}\label{eq:main}
  \Bigl|\int_{[\gamma,\infty)}\Phi(\sqrt\alpha E)\,d\rho(E)\Bigr|
  \le \Bigl|\int_{(-\infty,-\gamma]}\Phi(\sqrt\alpha E)\,d\tilde\sigma(E)\Bigr|
      +2M_0\frac{\alpha}{\pi R^2}e^{-R^2/2\alpha}.
\end{equation}
Now use the Gaussian tail bound $\Phi(-x)\le\tfrac12e^{-x^2/2}$ for $x\ge0$. On
$(-\infty,-\gamma]$ we have $\Phi(\sqrt\alpha E)\le\Phi(-\sqrt\alpha\gamma)\le\tfrac12e^{-\alpha\gamma^2/2}$,
so the first term on the right of \eqref{eq:main} is at most $\tfrac12M_0e^{-\alpha\gamma^2/2}$.
On $[\gamma,\infty)$ we have $|1-\Phi(\sqrt\alpha E)|\le\tfrac12e^{-\alpha\gamma^2/2}$, so the left
side of \eqref{eq:main} is at least $|\rho([\gamma,\infty))|-\tfrac12M_0e^{-\alpha\gamma^2/2}
=|F(0)|-\tfrac12M_0e^{-\alpha\gamma^2/2}$. Therefore
\[
  |F(0)|\ \le\ M_0\Bigl(e^{-\alpha\gamma^2/2}+\frac{2\alpha}{\pi R^2}e^{-R^2/2\alpha}\Bigr).
\]
Choosing $\alpha=R/\gamma$ equalises the exponents at $\gamma R/2$ and yields
$|F(0)|\le M_0e^{-\gamma R/2}\bigl(1+\tfrac{2}{\pi\gamma R}\bigr)\le2M_0e^{-\gamma R/2}$ when
$\gamma R\ge1$. Lemma~\ref{lem:FG}(i) finishes the proof.
\end{proof}

\begin{remark}
Nothing in the proof is special to $\omega_g$: replacing $(\Om_g,H(g))$ by $(\Om_\omega,H_\omega)$
gives the same bound for any ground state $\omega$ of $(\A,\alpha)$ whose GNS Hamiltonian has a
gap $\gamma$ above $0$. Note also where the sharp light cone is used --- Lemma~\ref{lem:FG}(iii)
holds \emph{exactly} for $|t|<R$. With only a Lieb--Robinson bound one would pay an extra tail
here, and the rate $\gamma/2$ would degrade to $\gamma/2v$ with additional error terms.
\end{remark}

\begin{corollary}\label{cor:unif}
Under Hypothesis~\ref{hyp:gap}, the bound of Theorem~\ref{thm:clust} holds for every
$\omega_{\tilde g_n}$ with the same $\gamma$, and hence for every limit point of
$(\omega_{\tilde g_n})$ in the topology of \cite[Cor.~4.3]{M0}.
\end{corollary}

\begin{remark}[the rate is off by a factor $2$, and this is intrinsic to the method]
\label{rem:factor2}
For the free field of mass $m$ one has $\gamma=m$, while the connected two-point function decays
as $e^{-mR}$ (the kernel of $(2\mu)^{-1}$ in one dimension behaves like $K_0(mR)$). So the true
rate is $\gamma$, whereas Theorem~\ref{thm:clust} delivers $\gamma/2$. The loss comes from the
optimisation $\alpha=R/\gamma$, which balances a Gaussian energy cutoff against a Gaussian
spatial tail and is the same factor-$2$ loss present in the Nachtergaele--Sims argument. It is
harmless for \S\ref{sec:dich}, where only $\gamma>0$ matters, but should be remembered when
Theorem~\ref{thm:clust} is fed into a rate computation: combined with M2, whose sharp exponent is
$2m_1$, one should not expect the two constants to compose naively.
\end{remark}

\section{Clustering does not give uniqueness}\label{sec:notuniq}

\begin{remark}[withdrawal of a claim]\label{rem:withdraw}
The strategy note, \S4.2, argued: uniform gap $\Rightarrow$ exponential clustering $\Rightarrow$
uniqueness ``by a standard ergodicity argument''. The first implication is
Theorem~\ref{thm:clust}. The second is false. What ergodicity gives is that a translation-invariant
clustering state is \emph{extremal} in the set of translation-invariant states
\cite[Thm.~4.3.17]{BR1}. Extremality is compatible with non-uniqueness: in a phase with broken
$\Z_2$ symmetry there are two extremal ground states $\omega_\pm$, and \emph{both} cluster
exponentially. Clustering separates pure states from mixtures; it does not count the pure states.
\end{remark}

What clustering does give is the following, which is genuinely useful because the limit points
produced by \cite{M0,M1} are $\Z_2$-invariant.

\section{The symmetry-breaking dichotomy}\label{sec:dich}

\begin{hypothesis}[structure of the broken phase]\label{hyp:broken}
$P$ is even, and there exist ground states $\omega_\pm$ of $(\A,\alpha)$ with
$\theta\omega_\pm=\omega_\mp$, both translation invariant and clustering, and an
$A=A^*\in\A(O)$ with $\omega_\pm(A)=\pm m$, $m\ne0$; moreover the $\Z_2$-invariant ground states
are exactly the segment $\{\lambda\omega_++(1-\lambda)\omega_-\}$ with $\lambda=\tfrac12$.
\end{hypothesis}

Hypothesis~\ref{hyp:broken} is what \cite{GJS76} is designed to establish at strong coupling.

\begin{theorem}[dichotomy]\label{thm:dich}
Assume $P$ even and Hypothesis~\ref{hyp:broken}. Then
$\inf_n\gapop H(\tilde g_n)=0$; that is, Hypothesis~\ref{hyp:gap} fails.
\end{theorem}

\begin{proof}
Suppose $\gamma:=\inf_n\gapop H(\tilde g_n)>0$. By \cite[Cor.~4.3]{M0} and \cite[Cor.~4.1]{M1}
there is a limit point $\omega$ of $(\omega_{\tilde g_n})$, a ground state of $(\A,\alpha)$, and by
\cite[Lem.~3.3]{M0} it is $\Z_2$-invariant. By Corollary~\ref{cor:unif}, $\omega$ satisfies the
bound of Theorem~\ref{thm:clust}, so for $A$ as in Hypothesis~\ref{hyp:broken},
\[
  \bigl|\omega\bigl(A\,\tau_a(A)\bigr)-\omega(A)^2\bigr|\ \le\ 2\|A\|^2e^{-\gamma\,\dist(O,O+a)/2}
  \ \xrightarrow[\ |a|\to\infty\ ]{}\ 0 .
\]
By Hypothesis~\ref{hyp:broken}, $\omega=\tfrac12(\omega_++\omega_-)$, so $\omega(A)=0$ while
\[
  \omega\bigl(A\tau_a(A)\bigr)=\tfrac12\Bigl[\omega_+\bigl(A\tau_a(A)\bigr)+\omega_-\bigl(A\tau_a(A)\bigr)\Bigr]
  \xrightarrow[\ |a|\to\infty\ ]{}\ \tfrac12\bigl[m^2+m^2\bigr]=m^2\ne0,
\]
using clustering and translation invariance of $\omega_\pm$. This contradicts the previous display.
\end{proof}

\begin{remark}
Theorem~\ref{thm:dich} is the expected physics --- in a symmetry-broken phase the finite-volume
gap closes (tunnelling splitting) --- but it is worth having as a theorem, for two reasons. First,
it shows Hypothesis~\ref{hyp:gap} is not a removable technicality: it is \emph{equivalent} to
absence of symmetry breaking along the sequence, so any proof of (H3) must be a proof that the
phase is unbroken. Second, it is a nontrivial internal consistency check on the M0--M2
architecture: the $\Z_2$-invariance of the limit points, which came for free from
$\Om_g>0$, is exactly what makes the contradiction work.
\end{remark}

\section{Reduction of uniqueness}\label{sec:red}

Write $\omega_\infty$ for the state obtained from the infinite-volume Euclidean measure $\nu$.

\begin{hypothesis}\label{hyp:R}\hfill
\begin{enumerate}
\item[\rm(R1)] \emph{(Euclidean uniqueness.)} The infinite-volume $\varphi^4_2$ measure $\nu$ is the
      unique translation-invariant solution of the relevant DLR / stochastic-quantisation
      characterisation, in the regime considered.
\item[\rm(R2)] \emph{(Global Markov property.)} $\nu$ has the global Markov property, equivalently
      the time-zero fields are cyclic, so that $\nu$ determines a state $\omega_\infty$ on the
      time-zero algebra $\A$.
\item[\rm(R3)] \emph{(Reverse reconstruction.)} Every locally normal, translation-invariant ground
      state $\omega$ of $(\A,\alpha)$ has Wightman functions that are tempered and whose analytic
      continuation yields Schwinger functions of a measure $\nu_\omega$ satisfying the
      characterisation in (R1).
\item[\rm(R4)] \emph{(Translation invariance of the limit.)} Every limit point of
      $(\omega_{\tilde g_n})$ is translation invariant.
\end{enumerate}
\end{hypothesis}

\begin{theorem}\label{thm:red}
Assume (R1)--(R4). Then $(\A,\alpha)$ has a unique locally normal translation-invariant ground
state, namely $\omega_\infty$, and $\omega_{\tilde g_n}\to\omega_\infty$ in norm on every $\A(O)$.
\end{theorem}

\begin{proof}
Let $\omega$ be a locally normal translation-invariant ground state. By (R3) it yields
$\nu_\omega$ satisfying the characterisation of (R1); by (R1), $\nu_\omega=\nu$; by (R2), $\nu$
determines a unique state on $\A$, so $\omega=\omega_\infty$. Thus the ground state is unique. By
\cite[Cor.~4.3]{M0} and \cite[Cor.~4.1]{M1} every limit point of $(\omega_{\tilde g_n})$ is a
locally normal ground state, and by (R4) it is translation invariant, hence equals
$\omega_\infty$; \cite[Cor.~4.2]{M1} upgrades this to convergence of the sequence.
\end{proof}

\begin{remark}[where the hypotheses come from, and which is open]\label{rem:status}
(R1) is available: Bauerschmidt--Dagallier--Weber (uniqueness of the invariant measure of the
$\varphi^4_2$ SPDE up to the critical temperature) and Albeverio--H\o egh-Krohn--Zegarlinski
(Prop.~7.2: uniqueness at differentiability points of the pressure in the external field).
(R2) is available at weak coupling: Albeverio--H\o egh-Krohn--Zegarlinski prove the global Markov
property for weak-coupling $\varphi^4_2$, and by Albeverio--H\o egh-Krohn the global Markov
property is equivalent to cyclicity of the time-zero fields.
(R4) would follow from uniqueness itself by \cite[Cor.~2.3]{M0}, so it is not an independent
assumption once (R1)--(R3) give uniqueness within the translation-invariant class; alternatively it
can be enforced by averaging, at the cost of \cite[Rem.~3.4]{M0}.

\textbf{(R3) is the open one.} It is a reconstruction in the direction opposite to
Osterwalder--Schrader: from an algebraic ground state to a Euclidean measure. Positivity of the
energy supplies the analytic continuation; what must be supplied in addition is temperedness
(for which \cite{GJIV} Thm~$A_g$ is the natural input) and the identification of the resulting
Schwinger functions as solutions of the correct Dyson--Schwinger hierarchy, which is where the
interaction enters. No proof of (R3) exists in the literature.
\end{remark}

\section{Status}

\begin{itemize}
\item \textbf{Proved here:} Theorem~\ref{thm:clust} (exponential clustering, absolute constant,
      rate $\gamma/2$, no Lieb--Robinson input); Theorem~\ref{thm:dich} (uniform gap fails in a
      symmetry-broken phase).
\item \textbf{Corrected here:} the uniqueness claim of the strategy note \S4.2 is withdrawn
      (Remark~\ref{rem:withdraw}).
\item \textbf{Reduced here:} uniqueness follows from (R1)--(R4), Theorem~\ref{thm:red}.
\item \textbf{Open:} (R3). \textbf{M3 is half-closed.} The purely Hamiltonian route of strategy
      \S4.2 is dead as stated; the Euclidean route of \S4.1 survives and now has a precise
      statement of its single missing ingredient.
\end{itemize}

\begin{thebibliography}{9}
\bibitem{M0} \emph{Limit points of the finite-volume ground states of $P(\varphi)_2$ are ground
states of the quasi-local dynamics}, Milestone M0, \texttt{../m0/m0.tex}.
\bibitem{M1} \emph{The local number bound without space-averaging}, Milestone M1,
\texttt{../m1/m1.tex}.
\bibitem{BR1} O.~Bratteli and D.~W.~Robinson, \emph{Operator Algebras and Quantum Statistical
Mechanics 1}, 2nd ed., Springer, 1987.
\bibitem{GJII} J.~Glimm and A.~Jaffe, \emph{The $\lambda(\phi^4)_2$ quantum field theory without
cutoffs. II}, Ann.\ of Math.\ \textbf{91} (1970) 362--401.
\bibitem{GJIV} J.~Glimm and A.~Jaffe, \emph{The $\lambda\phi^4_2$ quantum field theory without
cutoffs. IV}, J.\ Math.\ Phys.\ \textbf{13} (1972) 1568--1584.
\bibitem{GJS76} J.~Glimm, A.~Jaffe and T.~Spencer, \emph{A convergent expansion about mean field
theory I, II}, Ann.\ Physics \textbf{101} (1976) 610--630, 631--669.
\end{thebibliography}

\end{document}
